% Part: first-order-logic
% Chapter: natural-deduction
% Section: identity

\documentclass[../../../include/open-logic-section]{subfiles}

\begin{document}

\olfileid{fol}{ntd}{ide}

\olsection{\usetoken{P}{derivation} सह \usetoken{S}{identity}}

!!{identity} असलेल्या !!^{derivation}s साठी अतिरिक्त अनुमाननियम आवश्यक असतात.

\begin{defish}
\AxiomC{}
\RightLabel{\Intro{\eq}}
\UnaryInfC{$\eq[t][t]$}
\DisplayProof
\hfill
\begin{tabular}{r}
\AxiomC{$\eq[t_1][t_2]$}
\AxiomC{$!A(t_1)$}
\RightLabel{\Elim{\eq}}
\BinaryInfC{$!A(t_2)$}
\DisplayProof
\\[3ex]
\AxiomC{$\eq[t_1][t_2]$}
\AxiomC{$!A(t_2)$}
\RightLabel{\Elim{\eq}}
\BinaryInfC{$!A(t_1)$}
\DisplayProof
\end{tabular}
\end{defish}

वरील नियमांमध्ये $t$, $t_1$ आणि $t_2$ ही बंद पदे आहेत. \Intro{\eq} नियम
आपल्याला कोणत्याही गृहीतकांशिवाय, थेट $\eq[t][t]$ या रूपाचे कोणतेही
एकरूपता-विधान !!{derive} करू देतो.

\begin{ex}
$s$ आणि $t$ ही बंद पदे असतील, तर $!A(s), \eq[s][t] \Proves !A(t)$:
\begin{prooftree}
\AxiomC{$\eq[s][t]$}
\AxiomC{$!A(s)$}
\RightLabel{$\Elim{\eq}$}
\BinaryInfC{$!A(t)$}
\end{prooftree}
याला “एकरूप वस्तूंच्या प्रतिस्थापनाचे तत्त्व” किंवा लाइब्निझचा नियम
म्हणून ओळखले जाऊ शकते.
\end{ex}

\begin{prob}
$=$ हे सममित आणि संक्रमक दोन्ही आहे हे सिद्ध करा; म्हणजे
$\lforall[x][\lforall[y][(\eq[x][y] \lif
    \eq[y][x])]]$ आणि $\lforall[x][\lforall[y][\lforall[z]((\eq[x][y]
    \land \eq[y][z]) \lif \eq[x][z])]]$ यांच्या
!!{derivation}s द्या.
\end{prob}

\begin{ex}
आपण पुढील !!{sentence} !!{derive} करतो:
\begin{align*}
& \lforall[x][\lforall[y][((!A(x) \land !A(y)) \lif \eq[x][y])]]
\intertext{या !!{sentence} पासून:}
& \lexists[x][\lforall[y][(!A(y) \lif \eq[y][x])]]
\end{align*}
!!{derivation} उलट्या दिशेने विकसित करू:
\begin{prooftree}
\AxiomC{$\lexists[x][\lforall[y][(!A(y) \lif \eq[y][x])]]
\quad \Discharge{!A(a) \land !A(b)}{1}$}
\DeduceC{$\eq[a][b]$}
\DischargeRule{\Intro{\lif}}{1}
\UnaryInfC{$((!A(a) \land !A(b)) \lif \eq[a][b])$}
\RightLabel{\Intro{\lforall}}
\UnaryInfC{$\lforall[y][((!A(a) \land !A(y)) \lif \eq[a][y])]$}
\RightLabel{\Intro{\lforall}}
\UnaryInfC{$\lforall[x][\lforall[y][((!A(x) \land !A(y)) \lif \eq[x][y])]]$}
\end{prooftree}
आता मुख्य गृहीतक वापरावे लागेल: ते अस्तित्ववाची !!{formula} असल्यामुळे
\Elim{\lexists} वापरून मध्यवर्ती निष्कर्ष~$\eq[a][b]$ !!{derive} करू.
\begin{prooftree}
\AxiomC{$\lexists[x][\lforall[y][(!A(y) \lif \eq[y][x])]]$}
\AxiomC{$\Discharge{\lforall[y][(!A(y) \lif \eq[y][c]])}{2}$}
\noLine
\UnaryInfC{$\Discharge{!A(a) \land !A(b)}{1}$}
\DeduceC{$\eq[a][b]$}
\DischargeRule{\Elim{\lexists}}{2}
\BinaryInfC{$\eq[a][b]$}
\DischargeRule{\Intro{\lif}}{1}
\UnaryInfC{$((!A(a) \land !A(b)) \lif \eq[a][b])$}
\RightLabel{\Intro{\lforall}}
\UnaryInfC{$\lforall[y][((!A(a) \land !A(y)) \lif \eq[a][y])]$}
\RightLabel{\Intro{\lforall}}
\UnaryInfC{$\lforall[x][\lforall[y][((!A(x) \land !A(y)) \lif \eq[x][y])]]$}
\end{prooftree}
वरच्या उजव्या बाजूची उप-!!{derivation} तिच्या गृहीतकांचा वापर करून
$\eq[a][c]$ आणि $\eq[b][c]$ दाखवल्यावर पूर्ण होते. यासाठी दोन स्वतंत्र
!!{derivation}s आवश्यक आहेत. $\eq[a][c]$ साठीची !!{derivation}
पुढीलप्रमाणे आहे:
\begin{prooftree}
\AxiomC{$\Discharge{\lforall[y][(!A(y) \lif \eq[y][c]])}{2}$}
\RightLabel{\Elim{\lforall}}
\UnaryInfC{$!A(a) \lif \eq[a][c]$}
\AxiomC{$\Discharge{!A(a) \land !A(b)}{1}$}
\RightLabel{\Elim{\land}}
\UnaryInfC{$!A(a)$}
\RightLabel{\Elim{\lif}}
\BinaryInfC{$\eq[a][c]$}
\end{prooftree}
$\eq[a][c]$ आणि $\eq[b][c]$ यांच्यापासून \Elim{\eq} वापरून
$\eq[a][b]$ !!{derive} करतो.
\end{ex}

\begin{prob}
पुढील !!{formula}s च्या !!{derivation}s द्या:
\begin{enumerate}
\item $\lforall[x][\lforall[y][((\eq[x][y] \land !A(x)) \lif !A(y))]]$
\item $\lexists[x][!A(x)] \land \lforall[y][\lforall[z][((!A(y) \land
    !A(z)) \lif \eq[y][z])]] \lif \lexists[x][(!A(x) \land
  \lforall[y][(!A(y) \lif \eq[y][x])])]$
\end{enumerate}
\end{prob}

\end{document}
